% chapters/10/section-02.tex
\section{Dari BFS ke Dijkstra ke A*}
\label{sec:bfs-dijkstra-astar}

Judul bab ini menyebut tiga algoritma. Bagian ini menjelaskan mengapa yang
pertama tidak memadai, apa yang ditambahkan yang kedua, dan apa yang
diringankan oleh yang ketiga.

\subsection{Pencarian Melebar: Benar tetapi Tidak Menghiraukan Biaya}

Pencarian melebar (\istilah{breadth-first search}, BFS) memeriksa simpul
seturut jaraknya dari titik awal: mula-mula seluruh tetangga titik awal, lalu
tetangga dari tetangga, dan seterusnya. Cara itu menjamin jalur terpendek
apabila seluruh sisi berbiaya sama.

Kelemahannya muncul ketika biaya sisi tidak sama. Apabila melintasi pasir
memerlukan biaya dua kali lipat melintasi jalan biasa, BFS tidak
mengetahuinya, sebab ia hanya menghitung banyaknya langkah. Jalur dengan
langkah paling sedikit belum tentu jalur termurah.

\subsection{Dijkstra: Memeriksa Menurut Biaya}

Algoritma Dijkstra memperbaiki hal itu dengan memeriksa simpul menurut
\textbf{biaya sejauh ini} dari titik awal, bukan menurut banyaknya langkah.
Simpul dengan biaya terkecil selalu diperiksa lebih dahulu, sehingga jalur
yang ditemukan dijamin termurah, berapa pun biaya setiap sisinya.

Biaya sejauh ini dilambangkan $g$. Pada bab ini seluruh sisi berbiaya satu,
sehingga $g$ juga menyatakan banyaknya langkah dari titik awal.

\subsection{A*: Menambahkan Perkiraan ke Tujuan}

Dijkstra menyebar ke segala arah secara merata, sebab ia tidak mengetahui di
mana tujuannya berada. Akibatnya banyak simpul yang sebenarnya tidak mungkin
menuju tujuan ikut diperiksa.

Algoritma A* memperbaiki hal itu dengan menambahkan \textbf{perkiraan sisa
jarak} ke tujuan, dilambangkan $h$, pada biaya sejauh ini. Jumlah keduanya
dilambangkan $f$ dan dipakai sebagai urutan pemeriksaan.

\begin{equation}
\label{eq:f-gh}
f = g + h .
\end{equation}

Simpul yang memiliki $f$ terkecil diperiksa lebih dahulu. Karena $h$
menambahkan perkiraan sisa jarak, simpul yang menuju ke arah tujuan memperoleh
$f$ yang lebih kecil daripada simpul yang menjauh, sehingga diperiksa lebih
dahulu. Pencarian menjadi terarah.

\begin{catatabox}
Dijkstra sesungguhnya adalah A* dengan $h$ bernilai nol.
Karena itu, keduanya dapat ditulis sebagai satu program dengan satu pilihan,
dan itulah yang dilakukan pada bab ini. Pilihan heuristik bernama
\code{tanpa} berarti $h = 0$, yaitu Dijkstra.
\end{catatabox}

\subsection{Antrean Prioritas}

Agar simpul dengan $f$ terkecil selalu diperiksa lebih dahulu, tepi pencarian
harus berupa antrean prioritas. Python menyediakannya pada modul
\code{heapq}.

\begin{kodekeluaran}
import heapq

tepi = [(heuristik(awal, tujuan), 0, awal)]
while tepi:
    _, biaya_sekarang, sekarang = heapq.heappop(tepi)
    for tetangga_petak in tetangga(lebar, tinggi, dinding, sekarang):
        biaya_baru = biaya_sekarang + 1
        if tetangga_petak not in biaya or biaya_baru < biaya[tetangga_petak]:
            biaya[tetangga_petak] = biaya_baru
            perkiraan = biaya_baru + heuristik(tetangga_petak, tujuan)
            heapq.heappush(tepi, (perkiraan, biaya_baru, tetangga_petak))
            asal[tetangga_petak] = sekarang
\end{kodekeluaran}

Perhatikan tiga hal pada potongan kode itu.

\begin{enumerate}
  \item \textbf{Petak terdepan selalu yang paling menjanjikan.}
        \texttt{heappop} mengeluarkan isi antrean dengan prioritas terkecil,
        yaitu petak dengan nilai $f$ terkecil.
  \item \textbf{Sebuah petak dapat masuk antrean lebih dari sekali.} Petak itu
        dimasukkan kembali setiap kali ditemukan jalur yang lebih murah
        menuju ke sana. Cara itu lebih sederhana daripada memperbarui isi
        antrean, dan hasilnya tetap benar.
  \item \textbf{Biaya yang sudah usang dilewati.} Pemeriksaan
        \code{biaya_sekarang > biaya.get(sekarang, 0)} membuang isi antrean
        yang tertinggal dari pencarian sebelumnya. Tanpa pemeriksaan itu,
        petak yang sama dapat diperiksa lebih dari sekali.
\end{enumerate}

\begin{tipsbox}
Isi antrean prioritas berupa rangkap tiga, yaitu
\code{(perkiraan, biaya_baru, tetangga_petak)}. Angka biaya diletakkan di
tengah dengan sengaja: apabila dua petak memiliki perkiraan yang sama, Python
membandingkan unsur berikutnya. Tanpa angka pembanding itu, Python akan
membandingkan petaknya, yaitu pasangan angka, dan itu berhasil juga---tetapi
urutan pemeriksaan menjadi tidak menentu sehingga dua kali menjalankan
program dapat menghasilkan jumlah pemeriksaan yang sedikit berbeda.
\end{tipsbox}
